% !TeX root = ../../Book2.tex
% 纲要标题：## 第23章*　配边、Pontryagin–Thom 与光滑结构
% 纲要位置：几何学计划纲要.md，第 1883 行。
% 正文状态：初稿；尚未正式审核。

% 纲要内容开始
% 
% > 本章目标与证明范围：以完整的带标架 Pontryagin–Thom 对应作为代表深结果，而不只画坍缩示意图。一般手术分类和 h-配边的完整证明仍外引，维数和范畴条件逐条声明。
% 
% 纲要内容结束

\OptionalChapter{配边、Pontryagin–Thom 与光滑结构}{g2:ch:23}
\BookChapterNumber{g2:ch:23}

\begin{chapterabstract}
	闭流形可以作为高一维紧流形的边界出现，配边由此把流形组成可计算的群。本章先证明 Stiefel--Whitney 数与 Pontryagin 数的配边不变性，再完整建立带标架 Pontryagin--Thom 对应：管状邻域坍缩与正则逆像互逆，增加环境坐标对应球面映射的悬挂。随后把手术解释为配边中的把手附着，说明基本群、中维交叉与 Whitney 技巧怎样进入分类。最后比较光滑、PL 与拓扑结构，用异光滑球面和四维 Euclid 空间的精确定理展示范畴差异。
\end{chapterabstract}

\begin{learninggoals}
	\begin{itemize}
		\item 区分无定向、定向及带稳定标架的配边，利用领圈构造群运算并核对边界符号。
		\item 用相对基本类证明示性数的配边不变性，并用射影空间检验流形能否成为边界。
		\item 完成坍缩映射与正则逆像的双向构造，分别处理管状邻域、标架及环境嵌入的选择。
		\item 把增加法方向与悬挂对应起来，证明 \(\Omega_n^{\mathrm{fr}}\cong\pi_n^{\mathrm S}\)。
		\item 读出手术迹的把手指标及其基本群效应，识别高维配边定理与四维光滑问题的适用条件。
	\end{itemize}
\end{learninggoals}

在球面上把一条纬线逐渐推到北极，纬线会缩成一点，随后消失。若要求每个时刻都是同一张圆周的嵌入，这个过程在消失时失败；若把所有时刻放在一起，轨迹却是一张光滑圆盘。这提示我们改变比较对象的方法：允许中间截面的拓扑发生变化，要求全部轨迹组成一张高一维的光滑流形。配边正是这一想法的整体形式。

\ProseParagraphBreak

前两章已经把切丛变成上同调类，再通过基本类求值得到整数或模二数。为什么某些示性数能阻止一个流形成为边界？在配边中，边界的稳定切丛来自内部，而整个边界的基本类在内部同调于零。这两条事实便给出不变性的证明。再多保留一组稳定标架，流形的信息还能转换成球面之间的映射；第9章引入的稳定同伦群于是获得了直接的几何解释。

\ProseParagraphBreak

本章的主要证明是带标架 Pontryagin--Thom 对应。所用横截性、管状邻域、领圈和固有淹没分别回到第14、15章；稳定化使用第9章第5节的悬挂定理。后两节从第17章的把手附着转入手术及分类问题。\(h\)-配边、手术障碍、结构提升和异光滑结构的深层结论将在陈述处明确引用正式文献，维数、边界及范畴均写入假设，供读者在具体例子中判断能否使用。

% !TeX root = ../../Book2.tex
% 纲要标题：### 23.1 配边及其几何构造
% 纲要位置：几何学计划纲要.md，第 1887 行。
% 正文状态：初稿；尚未正式审核。

% 纲要内容开始
% 
% * 定向、无定向和带标架配边
% * 粘合、逆元与群结构
% * 示性数的配边不变性
% 
% 纲要内容结束

\section{配边及其几何构造}
\label{g2:sec:2301}

球面是球的边界，环面是实心环面的边界。这两张曲面的拓扑不同，却都能由一张高一维的紧流形填满。若把两个闭流形一起看成某张紧流形的全部边界，就得到一种允许拓扑改变的比较关系。需要同时决定的是：填充过程中应保留哪些几何资料？只要求光滑结构、再要求方向、进一步指定切丛的稳定标架，会得到不同的配边关系。

\subsection{定向、无定向和带标架配边}
\label{g2:subsec:2301:01}
本节的闭流形均指紧且无边界的光滑流形，可以不连通。配边自身允许有边界，但也要求紧；例如用无限延伸的圆柱填充一个端面，不能作为这里的配边。

\begin{definition}\label{g2:def:unoriented-oriented-bordism}
	设 \(M_0,M_1\) 为闭光滑 \(n\) 维流形。若存在紧光滑 \(n+1\) 维流形 \(W\)，使
	\[
		\partial W\cong M_0\amalg M_1,
	\]
	就称 \(M_0\) 与 \(M_1\)\term[peibian]{配边}[bordant]；此时没有指定方向，亦称无定向配边。若 \(M_0,M_1,W\) 均定向，并且上述识别保向地写成
	\[
		\partial W\cong(-M_0)\amalg M_1,
	\]
	就称二者定向配边。这里 \(-M_0\) 表示反转方向，边界方向按外法向优先约定确定。把 \(M_0\) 称为入边界、\(M_1\) 称为出边界。若 \(M\) 与空流形配边，就称 \(M\) 为相应意义下的边界。
\end{definition}

外法向优先的具体含义是：若 \(\nu\) 指向 \(W\) 外部，则 \((v_1,\ldots,v_n)\) 在边界上为正向，恰好当 \((\nu,v_1,\ldots,v_n)\) 在 \(W\) 上为正向。区间的两个端点因而具有相反符号；一般的圆柱也同样有一端带负号。这与\cref{g2:prop:boundary-fundamental-class}中连接同态满足 \(\partial[W,\partial W]=[\partial W]\) 的约定一致。

例如 \(S^n=\partial D^{n+1}\) 是无定向边界，也是取诱导方向后的定向边界。\(\partial(D^2\times S^1)=S^1\times S^1\)，所以二维环面也为边界。将填充球面与环面的三维流形取不交并，就能得到球面到环面的定向配边。配边只记录两端是否能共同出现，\(W\) 内部没有要求每个中间截面都与两端微分同胚。

现在给切向量规定一组随点光滑变化的坐标。实际切丛可能没有全局标架，但加入若干平凡方向后可能出现标架；球面的径向补充就是例子。用 \(\varepsilon_M^r=M\times\mathbb R^r\) 表示带标准基的实平凡丛，底空间清楚时省略下标。

\begin{definition}\label{g2:def:stable-framed-bordism}
	闭光滑 \(n\) 维流形 \(M\) 上的一个\term[wendingqiebiaojia]{稳定切标架}[stable tangent framing]是丛同构
	\[
		\varphi:\mathrm TM\oplus\varepsilon^r\xrightarrow{\ \cong\ }\varepsilon^{n+r},
	\]
	其中 \(r\ge0\)。把 \(\varphi\) 换成 \(\varphi\oplus\operatorname{id}_{\varepsilon^s}\) 称为稳定化；两个稳定切标架在稳定到同一秩后若可经丛同构同伦相连，就视为同一稳定标架。带稳定切标架的流形之间的微分同胚，也要求保持此稳定同伦类。

	紧光滑带边界流形 \(W^{n+1}\) 的稳定切标架诱导边界标架，采用下面的外法向约定。若 \(\partial W\) 的标架分别等于 \(-\varphi_0\) 与 \(\varphi_1\)，就称 \((M_0,\varphi_0)\) 与 \((M_1,\varphi_1)\)\term[daibiaojiapeibian]{带标架配边}[framed bordant]。其中 \(-\varphi\) 是先稳定化，再将一个平凡方向的坐标反号所得的标架。
\end{definition}

把稳定方向的标准正向与目标空间的标准正向固定下来，\(\varphi\) 就确定了 \(\mathrm TM\) 的方向。但是方向只区分基变换行列式的正负，标架还指定了整组坐标如何随点变化；即使方向相同，两组标架之间也可能有无法消去的扭转。忘掉标架得到定向，忘掉定向得到无定向流形。边界上需要延拓的资料随之逐层减少。

下面把边界标架的符号写清楚，以便随后把带标架流形变成球面之间的映射。设
\[
	\Phi:\mathrm TW\oplus\varepsilon^r\longrightarrow\varepsilon^{n+1+r},
\]
取沿边界的向外横截向量场 \(\nu\)，令 \(v\in\mathrm T\partial W\)、\(a\in\mathbb R^r\)、\(t\in\mathbb R\)。将新增的稳定坐标放在最后，定义
\begin{equation}\label{g2:eq:boundary-tangent-framing}
	(\partial\Phi)(v,a,t)=\Phi\bigl(v+(-1)^{n+r}t\nu,a\bigr).
\end{equation}
从 \((v_1,\ldots,v_n,a_1,\ldots,a_r,\nu)\) 把 \(\nu\) 移到最前，经过 \(n+r\) 次交换，故式中的符号恰好保证诱导方向是外法向优先。公式中多出的一维正是边界在 \(W\) 内的法线，它解释了为何这里自然使用稳定标架。

这个定义不依赖所取向外横截向量场：这些向量场的正的法向分量可作线性插值，过程中始终横截，故所得标架同伦。若先把 \(\Phi\) 稳定化一次，式\eqref{g2:eq:boundary-tangent-framing}中的符号会改变，但新增的两个平凡坐标同时改变排列；两种边界标架相差末两坐标上的旋转 \((b,t)\mapsto(-t,b)\)，它经旋转角从 \(0\) 连续变到 \(\pi/2\) 与恒等相连。因此定义与稳定次数相容。类似地，反射哪一个平凡坐标也不影响逆标架的稳定同伦类：在两个平凡方向组成的平面内转动反射轴即可互相连接。

球面的径向同构
\[
	\mathrm TS^n\oplus\varepsilon^1\longrightarrow\varepsilon^{n+1},\qquad(v,t)\longmapsto v+tx
\]
在\cref{g2:eq:sphere-stable-tangent-characteristic}已经出现。按本节的边界约定，球内标准切标架诱导的是 \((v,t)\mapsto v+(-1)^ntx\)。这给出了每个球面上一组确实能延拓到球内的稳定标架。讨论带标架配边时，球面与它所选的标架一起构成对象。

\subsection{粘合、逆元与群结构}
\label{g2:subsec:2301:02}
若 \(W_{01}\) 从 \(M_0\) 配边到 \(M_1\)，\(W_{12}\) 从 \(M_1\) 配边到 \(M_2\)，应能沿共同的 \(M_1\) 把二者接起来。\cref{g2:thm:collar-neighborhood}提供了所需的光滑接缝：每个边界分支附近有 \(M_1\times[0,\varepsilon)\) 的坐标。将两侧的非负参数分别写成 \(u\) 和 \(-u\)，接缝邻域就成为 \(M_1\times(-\varepsilon,\varepsilon)\)。这在商空间上给出光滑流形结构，其剩余边界为 \((-M_0)\amalg M_2\)。定向情形在接缝处一侧为出边界、一侧为入边界，恰好使两个方向相接。

带标架情形还要把标架接起来。在领圈内利用乘积分解识别各切丛，标架就成为一族可逆矩阵；将领圈参数在靠近边界处平滑地压到零，便可把这族标架改成与参数无关。如果给定的边界标架只在稳定后同伦，先稳定到共同秩，再把该同伦放进一段圆柱，使两侧在接缝处逐点一致。对入边界反转外法向时，式\eqref{g2:eq:boundary-tangent-framing}的最后一坐标反号，正与入边界上的 \(-\varphi\) 相抵消。因此拼接后的 \(\mathrm TW\) 上得到同一组稳定标架。

圆柱同时说明了为何反转方向或反射标架应当给出逆元。若 \(\varphi:\mathrm TM\oplus\varepsilon^r\to\varepsilon^{n+r}\)，在 \(M\times[0,1]\) 上定义
\begin{equation}\label{g2:eq:framed-cylinder}
	\Psi(v+s\partial_u,a)=\bigl(\varphi(v,a),(-1)^{n+r}s\bigr).
\end{equation}
其稳定次数为 \(r\)。在上端 \(u=1\)，外法向为 \(\partial_u\)，边界公式给出 \((\varphi(v,a),t)\)；在下端 \(u=0\)，外法向为 \(-\partial_u\)，得到 \((\varphi(v,a),-t)\)。因此它的边界是 \((M,\varphi)\amalg(M,-\varphi)\)。忘掉标架，便恢复 \(\partial(M\times[0,1])=(-M)\amalg M\) 的定向圆柱。

\begin{proposition}\label{g2:prop:bordism-group}
	对每个整数 \(n\ge0\)，闭光滑 \(n\) 维流形的无定向配边、定向配边、带稳定切标架的配边都是等价关系。不交并诱导交换群，分别记作
	\[
		\Omega_n^{\mathrm O},\qquad\Omega_n^{\mathrm{SO}},\qquad\Omega_n^{\mathrm{fr}}.
	\]
	单位元是空流形的类。定向群中 \([M]\) 的逆元为 \([-M]\)，带标架群中 \([M,\varphi]\) 的逆元为 \([M,-\varphi]\)；无定向群中每个元素都是自身的逆元。忘掉标架、忘掉定向分别给出群同态
	\[
		\Omega_n^{\mathrm{fr}}\longrightarrow\Omega_n^{\mathrm{SO}}\longrightarrow\Omega_n^{\mathrm O}.
	\]
\end{proposition}
\begin{proof}
	圆柱给出自反性。交换入边界与出边界，定向情形同时反转 \(W\) 的方向，带标架情形同时反射 \(W\) 的一个稳定坐标，便得到反向配边；无定向情形只须交换边界标识。这证明对称性。上述领圈拼接证明传递性。

	若分别给出 \(M\) 到 \(M'\)、\(N\) 到 \(N'\) 的配边，取它们的不交并就得到 \(M\amalg N\) 到 \(M'\amalg N'\) 的配边，故不交并在等价类上良定义。不交并的结合与交换微分同胚保持各分支所带的结构，给出结合律与交换律，空流形给出单位元。圆柱的全部边界是 \((-M)\amalg M\)，或 \((M,-\varphi)\amalg(M,\varphi)\)，所以这些不交并代表零。无定向时同一个圆柱的边界是两份 \(M\)，故 \(2[M]=0\)。最后，保留较少资料仍是配边，而且与不交并相容，因此两个忘却映射都是群同态。
\end{proof}

\begin{example}\label{g2:exm:zero-dimensional-bordism}
	零维闭流形是有限个点。定向时每个点有一个符号，定义其代数点数为正点数减去负点数。若它是某个紧定向一维流形 \(W\) 的边界，则把常值上同调类 \(1\in H^0(W;\mathbb Z)\) 与边界基本类配对，有
	\[
		\langle1,[\partial W]\rangle
		=\langle1,i_*\partial[W,\partial W]\rangle=0,
	\]
	因为对的长正合列中 \(i_*\partial=0\)。所以代数点数在定向配边下不变。反之，代数点数为零时，可以将每个负点与一个正点配成一对，用一个定向区间连接；这些区间的不交并填满所有点。因此
	\[
		\Omega_0^{\mathrm{SO}}\cong\mathbb Z.
	\]
	无定向时改用 \(\mathbb F_2\) 系数，同一论证说明边界的点数必为偶数；偶数个点又可两两用区间配对，故 \(\Omega_0^{\mathrm O}\cong\mathbb F_2\)。这两个计算中，代数点数与奇偶数分别是完整的不变量。
\end{example}

\subsection{示性数的配边不变性}
\label{g2:subsec:2301:03}
第21章的示性类依赖切丛，第6章的基本类则允许对顶次上同调类求值。二者结合所得的 Stiefel--Whitney 数与 Pontryagin 数已在\cref{g2:def:characteristic-numbers}定义。配边提供了一条检验这些数的办法：若流形确是边界，它的稳定切丛来自填充流形，因而有关示性类也能延拓；能延拓的顶次类在整个边界上的求值为零。

\begin{theorem}\label{g2:thm:characteristic-number-bordism}
	设 \(M_0,M_1\) 为闭光滑 \(n\) 维流形。若二者无定向配边，则每个 Stiefel--Whitney 数相同：对正整数 \(i_1,\ldots,i_s\) 满足 \(\sum i_j=n\)，有
	\[
		\left\langle\prod_jw_{i_j}(\mathrm TM_0),[M_0]_2\right\rangle
		=\left\langle\prod_jw_{i_j}(\mathrm TM_1),[M_1]_2\right\rangle.
	\]
	若二者定向配边且 \(n=4k\)，则每个整数 Pontryagin 数相同：对 \(\sum i_j=k\)，有
	\[
		\left\langle\prod_jp_{i_j}(\mathrm TM_0),[M_0]\right\rangle
		=\left\langle\prod_jp_{i_j}(\mathrm TM_1),[M_1]\right\rangle.
	\]
	因此相应边界的上述各数均为零。每个 Stiefel--Whitney 数定义同态 \(\Omega_n^{\mathrm O}\to\mathbb F_2\)，每个 Pontryagin 数定义同态 \(\Omega_{4k}^{\mathrm{SO}}\to\mathbb Z\)。
\end{theorem}
\begin{proof}
	先证明有关切丛关系，再把求值转移到配边内。设 \(W\) 为相应的紧光滑配边，记 \(A=\partial W\)、\(i:A\hookrightarrow W\)。向外横截向量场给出实际丛同构
	\[
		\mathrm TA\oplus\varepsilon^1\cong i^*\mathrm TW.
	\]
	Stiefel--Whitney 类的 Whitney 公式及自然性（\cref{g2:thm:whitney-sum}）于是给出 \(w_j(\mathrm TA)=i^*w_j(\mathrm TW)\)。对于 Pontryagin 类，复化这一稳定同构，再用 \(p_j(E)=(-1)^jc_{2j}(E_{\mathbb C})\)，同样得到整数等式 \(p_j(\mathrm TA)=i^*p_j(\mathrm TW)\)。这里加的是平凡丛，其总陈类为 \(1\)，所以整数系数的稳定性直接成立。

	无定向情形令 \(R=\mathbb F_2\)，定向情形令 \(R=\mathbb Z\)。相对基本类与边界基本类满足
	\[
		[W,A]_R\in H_{n+1}(W,A;R),\qquad
		\partial[W,A]_R=[A]_R,
	\]
	这是\cref{g2:prop:relative-fundamental-class,g2:prop:boundary-fundamental-class}；模二系数时每个流形都具有该基本类。对的长正合列包含
	\[
		H_{n+1}(W,A;R)\xrightarrow{\partial}H_n(A;R)\xrightarrow{i_*}H_n(W;R),
	\]
	故 \(i_*[A]_R=0\)。对任意 \(\alpha\in H^n(W;R)\)，上同调与同调配对的自然性给出
	\begin{equation}\label{g2:eq:extendible-class-boundary-zero}
		\langle i^*\alpha,[A]_R\rangle=\langle\alpha,i_*[A]_R\rangle=0.
	\end{equation}
	取 \(\alpha=\prod_jw_{i_j}(\mathrm TW)\) 或 \(\alpha=\prod_jp_{i_j}(\mathrm TW)\)，前述切丛关系说明其边界限制恰是所需的示性类乘积。模二时 \([A]_2=[M_0]_2+[M_1]_2\)，而定向整数情形 \([A]=-[M_0]+[M_1]\)，所以零求值分别给出定理中的两条等式。

	在不交并上，基本类是各分支基本类之和，示性类逐分支限制为原来的示性类，故每个求值具有可加性。这也证明它们在配边群上诱导所述同态。
\end{proof}

这个证明的核心是式\eqref{g2:eq:extendible-class-boundary-zero}：全体边界在 \(W\) 内同调于零。它直接保留了模二与整数信息。Euler 数则还依赖实际切丛的秩，不能由 \(\mathrm TA\oplus\varepsilon^1\) 恢复；例如 \(S^2\) 是 \(D^3\) 的边界，却有 \(\chi(S^2)=2\)。

\begin{example}\label{g2:exm:projective-bordism-obstructions}
	在 \(\mathbb RP^2\) 上，令 \(a\in H^1(\mathbb RP^2;\mathbb F_2)\) 为非零类。式\eqref{g2:eq:real-projective-tangent-classes}给出
	\[
		w(\mathrm T\mathbb RP^2)=(1+a)^3=1+a+a^2.
	\]
	由于 \(\langle a^2,[\mathbb RP^2]_2\rangle=1\)，其 \(w_1^2\) 数与 \(w_2\) 数均为 \(1\)。因此 \(\mathbb RP^2\) 不能成为任何紧光滑三维流形的全部边界。

	给 \(\mathbb CP^2\) 取复定向，并令 \(h=c_1(\lambda^*)\) 满足 \(\langle h^2,[\mathbb CP^2]\rangle=1\)。由式\eqref{g2:eq:complex-projective-pontryagin}，
	\[
		p_1(\mathrm T\mathbb CP^2)=3h^2,\qquad
		\langle p_1,[\mathbb CP^2]\rangle=3.
	\]
	所以 \(\mathbb CP^2\) 不是定向边界。更进一步，\(m[\mathbb CP^2]\) 的此项 Pontryagin 数为 \(3m\)，当 \(m\ne0\) 时不会消失，故 \([\mathbb CP^2]\) 在 \(\Omega_4^{\mathrm{SO}}\) 中具有无限阶。
\end{example}

如果 \(M\) 带稳定切标架，\(\mathrm TM\oplus\varepsilon^r\) 平凡便使所有正次 Stiefel--Whitney 类与 Pontryagin 类为零。因而这些数可以先检验一张流形有没有稳定标架；但在已有标架后，它们不能分辨标架本身的扭转。下一节将用球面映射记录这部分资料。

陈数所需的结构还包括复结构。若 \(M^{2m}\) 的 \(\mathrm TM\oplus\varepsilon^{2r}\) 已选复结构，并且配边 \(W^{2m+1}\) 的稳定切丛 \(\mathrm TW\oplus\varepsilon^{2r+1}\) 也有复结构，按外法向约定在边界上诱导给定的稳定复结构，就可以讨论稳定复配边。此时边界陈类是 \(W\) 上陈类的限制，式\eqref{g2:eq:extendible-class-boundary-zero}的整数证明原样适用，得到陈数的配边不变性。普通定向配边只要求方向能延拓，故若要用陈数作障碍，还须指定这项复结构的延拓条件。
% 纲要条目排为编号小节标题，并进入目录与书签；说明段仍保留于纲要注释中。

% !TeX root = ../../Book2.tex
% 纲要标题：### 23.2 Pontryagin–Thom 对应
% 纲要位置：几何学计划纲要.md，第 1893 行。
% 正文状态：未撰写；本文件目前是框架。

% 纲要内容开始
% 
% * 嵌入和法丛标架
% * 管状邻域坍缩映射
% * 横截逆像构造反向映射
% * 同伦与配边的双向对应
% * 稳定维数与选择无关性
% 
%   将嵌入维数的增大与 II.9.5 的悬挂映射对应；分别证明管状邻域选择、标架同伦及环境嵌入选择不会改变稳定类。
% 
% 纲要内容结束

\section{Pontryagin–Thom 对应}
\label{g2:sec:2302}

待撰写

\subsection{嵌入和法丛标架}
\label{g2:subsec:2302:01}
待撰写

\subsection{管状邻域坍缩映射}
\label{g2:subsec:2302:02}
待撰写

\subsection{横截逆像构造反向映射}
\label{g2:subsec:2302:03}
待撰写

\subsection{同伦与配边的双向对应}
\label{g2:subsec:2302:04}
待撰写

\subsection{稳定维数与选择无关性}
\label{g2:subsec:2302:05}
待撰写
% 纲要条目排为编号小节标题，并进入目录与书签；说明段仍保留于纲要注释中。

% !TeX root = ../../Book2.tex
% 纲要标题：### 23.3 手术与把手
% 纲要位置：几何学计划纲要.md，第 1903 行。
% 正文状态：初稿；尚未正式审核。

% 纲要内容开始
% 
% * 局部手术和配边的把手分解
% * 基本群与中维障碍
% * Whitney 技巧
% * h-配边、s-配边与手术障碍
% 
% 纲要内容结束

\section{手术与把手}
\label{g2:sec:2303}

配边把两个流形放在同一个高一维流形的边界上。若沿高度切开这个配边，中间截面会怎样变化？第17章的把手附着已经给出答案所需的局部模型。本节从这个模型出发，说明怎样有目的地改变基本群，以及为什么同样的办法走到中维时会遇到相交问题。

\subsection{局部手术和配边的把手分解}
\label{g2:subsec:2303:01}
设 \(M\) 是闭光滑 \(n\) 维流形，\(0\le p\le n-1\)。要在一个 \(p\) 维球面附近动手，首先需要的不只是嵌入 \(S^p\hookrightarrow M\)，而是一个乘积邻域
\[
	\varphi:S^p\times D^{n-p}\hookrightarrow M.
\]
它等价于给中心球面选择一组法标架，再应用管状邻域定理。记 \(C=M\setminus\varphi(S^p\times\operatorname{int}D^{n-p})\)。两块
\[
	S^p\times D^{n-p},\qquad D^{p+1}\times S^{n-p-1}
\]
具有同一个边界 \(S^p\times S^{n-p-1}\)，因而可以相互替换。

\begin{definition}\label{g2:def:surgery}
	沿 \(\varphi\) 的\term[shoushu]{\(p\) 维手术}[surgery]是流形
	\[
		M_\varphi=C\cup_{\varphi|_{S^p\times S^{n-p-1}}}(D^{p+1}\times S^{n-p-1}),
	\]
	粘合处使用领圈并磨圆。把 \(D^{p+1}\times D^{n-p}\) 沿 \(S^p\times D^{n-p}\) 接到 \(M\times\{1\}\)，所得 \((n+1)\) 维配边称为这次手术的\term[shoushuji]{手术迹}[trace of surgery]，记为 \(W_\varphi\)。
\end{definition}

\begin{proposition}\label{g2:prop:surgery-trace}
	\(W_\varphi\) 从 \(M\) 出发只有一个指标 \(p+1\) 的把手，另一端边界是 \(M_\varphi\)。反向看，它从 \(M_\varphi\) 出发只有一个指标 \(n-p\) 的把手。若 \(M\) 定向，按定向兼容的附着方式构造，则 \(\partial W_\varphi=(-M)\amalg M_\varphi\)。
\end{proposition}
\begin{proof}
	乘积圆盘的边界由两个面组成：
	\[
		\partial(D^{p+1}\times D^{n-p})
		=(S^p\times D^{n-p})\cup(D^{p+1}\times S^{n-p-1}).
	\]
	第一个面附着后成为内部接缝，第二个面与 \(C\) 合成新边界。这正是定义中的替换。对照\cref{g2:def:morse-handle}，核心圆盘是 \(D^{p+1}\times\{0\}\)，故指标为 \(p+1\)。倒转配边方向，交换两个圆盘因子，原来的自由面成为附着面，新的核心为 \(D^{n-p}\)，指标成为 \(n-p\)。最后按外法向优先的边界定向约定，乘积的下端带相反定向，上端带原定向，得到所述边界公式。
\end{proof}

三个维数分别是：接受手术的流形维数 \(n\)、被取走的中心球面维数 \(p\)、手术迹的把手指标 \(p+1\)。这里把“\(p\) 维手术”按球面维数命名；其他文献还可能按把手指标命名。

\begin{proposition}\label{g2:prop:surgery-euler}
	一次 \(p\) 维手术满足
	\[
		\chi(M_\varphi)-\chi(M)=(-1)^{n-p-1}-(-1)^p.
	\]
	因此奇数维时 Euler 示性数不变，偶数维时改变量为 \(-2(-1)^p\)。
\end{proposition}
\begin{proof}
	利用领圈把粘合分解加厚成开覆盖，Mayer--Vietoris 序列给出 Euler 示性数的加法公式。两次分解具有相同的 \(C\) 和交界 \(S^p\times S^{n-p-1}\)，相减后只剩
	\[
		\chi(D^{p+1}\times S^{n-p-1})-\chi(S^p\times D^{n-p})
		=1+(-1)^{n-p-1}-1-(-1)^p.
	\]
	球面的 Euler 示性数来自它的零胞腔和顶维胞腔，圆盘可缩；故公式也包括 \(S^0\) 出现的端点情形。
\end{proof}

\begin{example}\label{g2:exm:surface-surgery}
	在定向曲面上，\(p=0\) 的手术取走两个小圆盘，再用圆柱 \(D^1\times S^1\) 接起两个边界圆。如果两个圆盘位于同一连通曲面，就增加一个把手，\(\chi\) 减少 \(2\)；如果它们分别位于两张曲面，就把两张曲面连成其连通和。反向的 \(p=1\) 手术切去一个圆柱邻域，再用两个圆盘封口，\(\chi\) 增加 \(2\)。沿不分离圆周切开会减少一个把手，沿分离圆周切开则产生两个分支。相同的局部替换可以有不同的全局效果，区别来自附着球面的位置。
\end{example}

任意紧光滑配边 \((W;M_0,M_1)\) 都可以这样逐步扫描。\Cref{g2:prop:boundary-morse-function}给出在两端领圈内等于高度参数、内部只有有限个非退化临界点且临界值互异的函数。逐次应用\cref{g2:thm:morse-handle-attachment}，每跨过一个指标 \(\lambda\) 的临界点，截面便经历一次 \((\lambda-1)\) 维手术。指标 \(0\) 和 \(n+1\) 分别使一个 \(S^n\) 分支诞生或消失，需要另计。这样得到的是有限把手分解；把手的位置和法标架保留了仅凭配边群元素看不到的信息。

\subsection{基本群与中维障碍}
\label{g2:subsec:2303:02}
先考察手术能否消去一个指定环路。直觉上，沿这个环路接上圆盘便会让它可缩；手术还必须同时补上横向厚度，保证结果仍是流形。手术迹将这两件事分开处理。

\begin{proposition}\label{g2:prop:surgery-fundamental-group}
	设 \(M\) 连通，\(1\le p\le n-3\)，并沿 \(\varphi:S^p\times D^{n-p}\hookrightarrow M\) 做手术。则 \(M_\varphi\) 连通，且
	\[
		\pi_1(M_\varphi)\cong
		\begin{cases}
			\pi_1(M)/\langle\!\langle[\alpha]\rangle\!\rangle,&p=1,\\
			\pi_1(M),&p\ge2.
		\end{cases}
	\]
	这里 \(\alpha=\varphi|_{S^1\times\{0\}}\)，用一条道路把它接到固定基点；\(\langle\!\langle[\alpha]\rangle\!\rangle\) 表示它生成的正规子群。更换这条道路只会共轭 \([\alpha]\)，不改变该正规子群。
\end{proposition}
\begin{proof}
	由\cref{g2:prop:surgery-trace,g2:lem:handle-core-retraction}，从 \(M\) 到 \(W_\varphi\) 的同伦效果是附着一个 \((p+1)\) 胞腔。对于 \(X\cup_\beta D^k\)，\(k\ge2\)，取圆盘内部及圆盘中心的补集作为开覆盖：前者可缩，后者收缩到 \(X\)，交集收缩到 \(S^{k-1}\)。\Cref{g2:thm:van-kampen}于是给出
	\[
		\pi_1(X\cup_\beta D^k)=
		\begin{cases}
			\pi_1(X)/\langle\!\langle\beta_*\pi_1(S^1)\rangle\!\rangle,&k=2,\\
			\pi_1(X),&k\ge3.
		\end{cases}
	\]
	第二行用到了 \(S^{k-1}\) 单连通。因而从 \(M\) 出发计算 \(\pi_1(W_\varphi)\)，正好得到命题右侧。

	反向的核心胞腔维数是 \(n-p\ge3\)。它的附着球面连通，不能把原有的不同连通分支接在一起；而 \(W_\varphi\) 从连通的 \(M\) 附着一个把手得到，故本身连通。因此 \(M_\varphi\) 也连通。对反向附着再用上式，\(\pi_1(M_\varphi)\to\pi_1(W_\varphi)\) 是同构。把两端的计算结合起来即得结论。
\end{proof}

维数限制在证明中有明确位置：\(n-p\ge3\) 保证反向附着不改变基本群。如果反向把手的指标降到 \(2\)，它也能消去关系；降到 \(1\)，则可能新增环路或连接分支。三维流形沿圆周的手术恰好落在前一种情形，因此不能套用上面的商群公式。

\ProseParagraphBreak

在高维定向流形中，有限组基本群生成元可用互不相交的嵌入圆周表示。每个圆周的法丛是定向的，而 \(S^1\) 上的定向实向量丛平凡，因而可以选择法标架。当 \(n\ge4\) 时，逐次作 \(1\) 维手术便能消去这些生成元。这说明存在同一有向配边类中的单连通代表；这一步改变了原流形的同伦型。

\ProseParagraphBreak

想继续消去更高同伦群，就要用嵌入的 \(S^p\) 代表待消去元素，并为它选择法标架。横截性提供一个有用的维数判断：两个 \(p\) 维代表在 \(n\) 维流形中的预期相交维数是 \(2p-n\)。在 \(2p<n\) 时，有希望通过一般位置使它们不相交；在 \(n=2p\) 时，交点成为孤立点，不能再靠一个小扰动全部移走。这是中维在手术理论中反复出现的原因。

\ProseParagraphBreak

设 \(M^{2q}\) 闭、连通、定向。其自由中维同调上的相交配对可写成
\[
	\lambda(x,y)=\bigl\langle\operatorname{PD}(x)\smile\operatorname{PD}(y),[M]\bigr\rangle,
	\qquad x,y\in H_q(M;\mathbb Z)/\text{挠子群}.
\]
这里 \(\operatorname{PD}(x)\in H^q(M;\mathbb Z)\) 是 Poincar\'e 对偶类，模挠意义下的取值不依赖代表；杯积的分次交换律给出 \(\lambda(y,x)=(-1)^q\lambda(x,y)\)。因而 \(q\) 偶数时得到对称形式，\(q\) 奇数时得到斜对称形式。若两个类由互不相交的球面代表，则交数为零。若一个球面具有平凡法丛，沿其中一个非零法向量场略作推移，就得到与自身不相交的同调代表，所以它的自交数也必须为零。法标架与相交数据都是能否实施手术的实际条件。

\begin{example}\label{g2:exm:surgery-hyperbolic-pair}
	在 \(S^q\times S^q\) 中，取 \(a=[S^q\times\{y\}]\)、\(b=[\{x\}\times S^q]\)，并选择乘积定向，则
	\[
		\bigl(\lambda(a_i,a_j)\bigr)_{(a_1,a_2)=(a,b)}
		=\begin{pmatrix}0&1\\(-1)^q&0\end{pmatrix}.
	\]
	两个因子球面各自都有平凡法丛，却具有非零交数，所以不可能同时由互不相交的代表实现。手术中的代数整理需要选择合适的半组生成元并保留与之对偶的数据；仅逐项检查每个生成元的自交数还不够。
\end{example}

基本群非平凡时，还需记录交点的道路信息。对两个带基点道路的、基本群像为零的球面代表，在交点 \(z\) 沿第一个球面到达 \(z\)，再沿第二个球面返回基点，得到 \(g_z\in\pi=\pi_1(M)\)。把交点符号 \(\varepsilon_z\) 一同记下，形成群环中的和
\[
	\lambda=\sum_z\varepsilon_z g_z\in\mathbb Z[\pi].
\]
固定基点道路后，这些标签可比较；改变选择会引入相应的群作用。取增广 \(\mathbb Z[\pi]\to\mathbb Z\)、把每个群元素送到 \(1\)，才得到普通整数交数。因此两个符号相反但群标签不同的交点，虽然在整数计数中相消，却未必能用下述圆盘消去。中维手术障碍保留群环配对及与自交、法标架有关的二次数据，原因就在这里。

\subsection{Whitney 技巧}
\label{g2:subsec:2303:03}
设定向子流形 \(A^a,B^b\subset X^{a+b}\) 横截相交，\(u,v\) 是符号相反的两个交点。在 \(A\) 中选一条连接 \(u,v\) 的嵌入弧，在 \(B\) 中再选一条反向连接的嵌入弧，使弧内部避开其他交点。它们合成一条带两个角的闭曲线，称为 Whitney 圆周。若能用一个圆盘填住它，就可以尝试把 \(A\) 的这一小片沿圆盘推过 \(B\)，使两次穿越一起消失。

\ProseParagraphBreak

这个圆盘还必须光滑嵌入，内部与 \(A\cup B\) 不相交，边界按上述两条弧落在两张片上。沿两条边界弧由片内方向给出的法向数据也须能延拓到整个圆盘，形成与两片相容的乘积邻域；这称为圆盘具有 Whitney 标架。只有闭曲线可缩，尚未检查这两件事。下面把已有这种局部乘积模型时的消点操作写出。

\begin{lemma}\label{g2:lem:local-whitney-move}
	设上述两个交点具有一个带 Whitney 标架的圆盘，其内部避开 \(A\cup B\)，且其小邻域内没有其他片。则存在支撑在该邻域中的光滑同痕，把 \(A\) 移成 \(A'\)，使
	\[
		A'\cap B=(A\cap B)\setminus\{u,v\}.
	\]
	这里同痕作用于 \(A\)，比较相交时保持 \(B\) 为原来的子流形。
\end{lemma}
\begin{proof}
	标架把圆盘的小邻域识别为一个平面邻域 \(U\) 乘上 \(\mathbb R^{a-1}\times\mathbb R^{b-1}\) 中的小圆盘。在这个识别下，两张片分别写成
	\[
		C_A\times\mathbb R^{a-1}\times\{0\},\qquad
		C_B\times\{0\}\times\mathbb R^{b-1},
	\]
	其中平面弧 \(C_A,C_B\) 围成待消去的二角形。把 \(C_A\) 推过这个二角形，可得到支撑在 \(U\) 内的平面同痕 \(G_t\)，使 \(G_1(C_A)\) 与 \(C_B\) 不相交。

	这个平面操作可以直接实现：把 \(C_B\) 拉直为横轴，把相关的 \(C_A\) 写成 \(y=h(x)\)，其中两个横截零点之间 \(h<0\)，两端外侧 \(h>0\)。选紧支光滑函数 \(k\ge0\)，使 \(h+k>0\) 于相关区间，令 \(h_t=h+tk\)。每条图像都嵌入；取沿图像等于 \(k(x)\partial_y\)、离开图像附近便截断为零的时变向量场，它的流将 \(h_0\) 送到 \(h_t\)，给出所需 \(G_t\)。二角形的两边嵌入且内部没有其他片，允许先作上述平面坐标调整。

	取横向变量上的光滑截断 \(\rho(x,y)\)，在 \((0,0)\) 附近为 \(1\)，靠近横向邻域边界为 \(0\)，定义
	\[
		H_t(z,x,y)=\bigl(G_{t\rho(x,y)}(z),x,y\bigr).
	\]
	每个固定 \((x,y)\) 的平面映射都是微分同胚，逆映射也光滑依赖参数，故 \(H_t\) 是同痕，并可与邻域外的恒等映射拼接。移动后的第一张片仍有 \(y=0\)，第二张片仍有 \(x=0\)，所以二者若相交，只能发生在 \(x=y=0\)。此处 \(\rho=1\)，相交等价于 \(G_1(C_A)\cap C_B\)，已经为空。邻域外没有变化，故恰好消去指定的两个交点。
\end{proof}

上述引理把几何操作与圆盘的存在问题分开了。以下是高维中常用的一组充分条件。

\begin{theorem}[Whitney 技巧]\label{g2:thm:whitney-trick}
	设 \(X^{a+b}\) 是单连通定向光滑流形，\(A^a,B^b\) 是其中连通、定向、横截相交的闭嵌入子流形。若 \(a,b\ge3\)，则任意两个符号相反的交点可以通过 \(A\) 的光滑同痕消去，而不改变其余交点。
\end{theorem}

这里引用完整的 Whitney 圆盘与标架构造，见 \cite[定理6.6及引理6.7]{Milnor1965HCobordism}；群环标签版本见 \cite[第7.3节]{Ranicki2002Surgery}。其维数机制如下：单连通性使 Whitney 圆周能够由圆盘填充；对填充圆盘作相对一般位置调整，可能出现的自交维数是 \(4-(a+b)<0\)，与 \(A\)、\(B\) 内部的预期相交维数分别是 \(2-b<0\)、\(2-a<0\)。这些不等式使圆盘可以嵌入且内部避开两片。剩下的边界法向扭转由标架构造处理，随后应用\cref{g2:lem:local-whitney-move}。标架这一步是引用定理的一部分，不能只由“圆盘可缩”替代。

\ProseParagraphBreak

在四维中若 \(a=b=2\)，刚才三个预期维数都变成 \(0\)。填充圆盘本身可以有双点，也可以穿过原来的两张曲面；试图消去一对交点可能又产生新的交点。单连通性仍解决边界圆周可缩的问题，却不能保证存在内部干净且带所需标架的嵌入圆盘。已经给定这种圆盘时，局部消点引理仍然成立；失去的是由一般位置自动得到圆盘的步骤。下一节的四维光滑现象与这道障碍密切相关。

\subsection{h-配边、s-配边与手术障碍}
\label{g2:subsec:2303:04}
配边允许中间截面改变许多次。若两端的包含映射都已经是同伦等价，这些把手在同伦论意义下没有留下新信息，能否把它们全部取消？

\begin{definition}\label{g2:def:h-cobordism}
	紧光滑配边 \((W;M_0,M_1)\) 称为\term[hpeibian]{\(h\)-配边}[h-cobordism]，如果两个包含 \(M_i\hookrightarrow W\) 都是同伦等价。称它为乘积配边，是指存在微分同胚
	\[
		F:M_0\times[0,1]\longrightarrow W,
		\qquad F(x,0)=x,\quad F(M_0\times\{1\})=M_1.
	\]
\end{definition}

\begin{theorem}[\(h\)-配边定理]\label{g2:thm:h-cobordism}
	设 \((W;M_0,M_1)\) 为光滑 \(h\)-配边，\(M_0,M_1\) 是闭连通流形，\(\dim W\ge6\)。如果 \(M_0\) 单连通，则 \(W\) 是乘积配边，特别地 \(M_0\) 与 \(M_1\) 微分同胚。
\end{theorem}

这是 Smale 的定理；完整证明可见 \cite[第6--9节，定理9.1]{Milnor1965HCobordism}。包含为同伦等价保证 \(W\) 和 \(M_1\) 也单连通。证明先把相对把手链复形中的代数消去实现为把手滑动，再用 Whitney 技巧把代数交数化成恰好一次的几何相交，最后取消相邻指标的把手。这个过程所用的高维把手整理是本节的外部输入。定理给出的是固定下端 \(M_0\) 的乘积结构；任意预先指定的上端参数化不一定也逐点固定。

\ProseParagraphBreak

非单连通时，相对万有覆盖的链复形是 \(\mathbb Z[\pi]\) 上的带基自由链复形，其中 \(\pi=\pi_1(M_0)\)。同伦等价使它链可缩，但把它约去所需的基变换未必能由初等变换及把基元素乘以 \(\pm g\)、\(g\in\pi\) 实现。这份差别由\term[whiteheadniuzi]{Whitehead 扭子}[Whitehead torsion]记录：
\[
	\tau(W,M_0)\in\operatorname{Wh}(\pi),\qquad
	\operatorname{Wh}(\pi)=K_1(\mathbb Z[\pi])/\langle[\pm g]:g\in\pi\rangle.
\]
这里 \(K_1(R)\) 是把可逆矩阵稳定化后再模去初等矩阵所得的群；利用链收缩构造奇、偶维链群之间的基变换，取其稳定类，便定义上述扭子。改变胞腔提升和定向会产生 \(\pm g\)，所以要再除掉这些类。构造及其不变性详见 \cite[第8.1--8.2节]{Ranicki2002Surgery}。

\begin{theorem}[\(s\)-配边定理]\label{g2:thm:s-cobordism}
	设 \((W;M_0,M_1)\) 为光滑 \(h\)-配边，两端闭且连通，\(\dim W\ge6\)。则它是固定 \(M_0\) 的乘积配边，当且仅当 \(\tau(W,M_0)=0\)。
\end{theorem}

这一版本见 \cite[第8.2节]{Ranicki2002Surgery}；扭子为零的 \(h\)-配边称为 \(s\)-配边。由于 \(\operatorname{Wh}(\{1\})=0\)，单连通版本是它的特例。使用定理时，首先要有两个包含映射的同伦等价，再检验扭子；两端同调群同构尚未给出这些条件。

\ProseParagraphBreak

最后说明手术怎样把一个“接近同伦等价”的映射改成同伦等价。设 \(M^n,N^n\) 闭、连通、定向。一个\term[duyifayingse]{度一法映射}[degree-one normal map]包含度为 \(1\) 的映射 \(f:M\to N\)、\(N\) 上的稳定实向量丛 \(\xi\)，以及稳定丛同构 \(b:\nu_M\cong f^*\xi\)。这里 \(\nu_M\) 是稳定法丛，同构也可表示为选定 \(\mathrm{T}M\oplus f^*\xi\) 的稳定平凡化。\(\xi\) 是数据的一部分，不预设为 \(\nu_N\)。称两份这样的数据\term[fapeibian]{法配边}[normally bordant]，是指存在配边上的映射及稳定法丛数据延拓它们，目标丛保持为 \(\xi\)。

\begin{theorem}[手术障碍定理]\label{g2:thm:simply-connected-surgery-obstruction}
	设 \(N^n\) 是单连通闭定向光滑流形，\(n\ge5\)，\((f,b):M^n\to N^n\) 是度一法映射。存在只依赖其法配边类的障碍
	\[
		\theta(f,b)\in L_n(\mathbb Z),
	\]
	它为零当且仅当 \((f,b)\) 法配边于一个同伦等价 \(f':M'\to N\)。这里
	\[
		L_n(\mathbb Z)\cong
		\begin{cases}
			\mathbb Z,&n\equiv0\pmod4,\\
			\mathbb Z/2,&n\equiv2\pmod4,\\
			0,&n\text{ 为奇数}.
		\end{cases}
	\]
\end{theorem}

这里引用手术障碍定理及单连通 \(L\) 群计算 \cite[定理3.2、13A.1]{Wall1999Surgery}。原定理表述为单同伦等价；此处基本群平凡，Whitehead 群为零，因而与同伦等价一致。\(L_n(\mathbb Z)\) 是二次手术障碍群：在偶数维，它保留中维相交形式及二次细化模去可通过标准手术消去的数据后的类。\(4k\) 维的整数由相交形式的符号差除以 \(8\) 给出，\(4k+2\) 维的二值障碍由相应中维二次形式的二值不变量给出；这些代数构造和上表的计算也包含在所引结果中。

\ProseParagraphBreak

例如，当目标维数为 \(5\) 或 \(7\) 时，一旦已经给出度一法映射，其障碍群为零，所以可以在法配边中得到同伦等价。输出是新的 \(M'\) 及 \(M'\simeq N\)；从同伦等价走到微分同胚还需进一步信息。对固定同伦类型的光滑流形进行分类，还须研究法数据之间的差别、手术群的作用和 \(s\)-配边关系。把这几个层次分清，下一节同一拓扑类型上不同光滑结构的例子才不会显得与手术理论矛盾。
% 纲要条目排为编号小节标题，并进入目录与书签；说明段仍保留于纲要注释中。

% !TeX root = ../../Book2.tex
% 纲要标题：### 23.4 光滑、PL 与拓扑类别
% 纲要位置：几何学计划纲要.md，第 1860 行。
% 正文状态：未撰写；本文件目前是框架。

% 纲要内容开始
% 
% * 结构的遗忘及提升问题
% * 低维和高维的不同现象
% * 典型异光滑结构的定理与限制
% 
% ---
% 
% 纲要内容结束

\section{光滑、PL 与拓扑类别}
\label{g2:sec:2304}

待撰写

\subsection{结构的遗忘及提升问题}
\label{g2:subsec:2304:01}
待撰写

\subsection{低维和高维的不同现象}
\label{g2:subsec:2304:02}
待撰写

\subsection{典型异光滑结构的定理与限制}
\label{g2:subsec:2304:03}
待撰写
% 纲要条目排为编号小节标题，并进入目录与书签；说明段仍保留于纲要注释中。


\begin{chaptersummary}
	配边把闭流形放到高一维紧流形的边界上比较。不交并给出群运算，圆柱给出逆元，领圈使配边能够拼接。边界稳定切丛来自内部，因此 Stiefel--Whitney 数和整数 Pontryagin 数在相应配边下不变；全体边界在内部同调于零是求值消失的关键。Euler 数还依赖实际秩，不能直接套用这一稳定性论证。

	\ProseParagraphBreak

	法标架把管状邻域的纤维统一成 \(\mathbb R^k\)，坍缩由此产生 \(S^{n+k}\to S^k\)。正则逆像通过微分取回子流形及其法标架。相对横截性把同伦变成配边，保持领圈的管状邻域又把配边变成同伦，两种构造因而互逆。增加一个环境坐标等于增加一个法标架向量，对应映射的悬挂；消去嵌入的选择后得到 \(\Omega_n^{\mathrm{fr}}\cong\pi_n^{\mathrm S}\)。

	\ProseParagraphBreak

	手术把 \(S^p\times D^{n-p}\) 换为 \(D^{p+1}\times S^{n-p-1}\)，其配边迹附着一个指标 \(p+1\) 的把手。低指标手术改变基本群，中维手术受到相交形式及标架数据的约束。高维的一般位置与 Whitney 消交支持 \(h\)-配边和手术理论；四维中圆盘与圆盘仍会相交，相关光滑问题呈现另一种结构。光滑、PL、拓扑三种等价关系须分别辨认，异光滑球面与异光滑 \(\mathbb R^4\) 正是同胚类型不足以确定光滑类型的具体实例。
\end{chaptersummary}

\BookExercises
% 第二卷第23章习题候选：对应23.1；由章入口统一收录。
\begin{exercise}\label{g2:ex:euler-boundary-formula}
	设 \(W\) 为紧光滑 \(d\) 维流形，\(d\ge1\)，不要求定向。用模二系数的 Poincaré--Lefschetz 对偶证明
	\[
		\chi(\partial W)=\bigl(1-(-1)^d\bigr)\chi(W).
	\]
	由此说明，Euler 示性数为奇数的闭光滑流形不能是无定向边界，并解释这与 \(S^2=\partial D^3\) 相容的原因。
\end{exercise}
\begin{solution}
	以下同调均取 \(\mathbb F_2\) 系数。先说明 Euler 交错和是有限的：把 \(W\) 双化为闭光滑流形 \(\operatorname{Dbl}(W)\)，由\cref{g2:cor:morse-cell-type}，双化具有有限 CW 同伦型。把两份副本都映回原来的 \(W\)，给出一个连续折叠映射，其在第一份 \(W\) 上是恒等，所以 \(H_k(W)\) 是双化的 \(H_k\) 的直和因子。边界本身是闭流形，其同调也有限；对的长正合列再保证相对同调有限。故这些 Euler 示性数均可由同调维数的有限交错和计算。

	\cref{g2:thm:poincare-lefschetz}在模二系数下给出
	\[
		\dim H_k(W,\partial W)=\dim H^{d-k}(W)=\dim H_{d-k}(W).
	\]
	因此 \(\chi(W,\partial W)=(-1)^d\chi(W)\)。对的长正合列中各项维数的交错和为零，于是
	\[
		\chi(W)=\chi(\partial W)+\chi(W,\partial W)
		=\chi(\partial W)+(-1)^d\chi(W),
	\]
	即得所求等式。若 \(d\) 偶，则边界的 Euler 示性数为零；若 \(d\) 奇，则它为 \(2\chi(W)\)。两种情形均为偶数，故奇数 Euler 示性数阻止流形成为无定向边界。对于 \(D^3\)，右端是 \(2\chi(D^3)=2\)，正好等于 \(\chi(S^2)\)；边界的 Euler 示性数可以非零。
\end{solution}

\begin{exercise}\label{g2:ex:projective-product-bordism-numbers}
	令 \(X=\mathbb RP^2\times\mathbb RP^2\)、\(Y=\mathbb CP^2\)。计算这两张实四维流形的全部 Stiefel--Whitney 数，即对应于
	\[
		w_1^4,\quad w_1^2w_2,\quad w_1w_3,\quad w_2^2,\quad w_4
	\]
	的五项求值。它们能否区分 \(X\) 与 \(Y\) 的无定向配边类？说明仅凭\cref{g2:thm:characteristic-number-bordism}能得到的结论。
\end{exercise}
\begin{solution}
	设 \(a,b\) 是从两因子拉回的一次生成元。模二上同调环为
	\[
		H^*(X;\mathbb F_2)=\mathbb F_2[a,b]/(a^3,b^3),\qquad
		\langle a^2b^2,[X]_2\rangle=1.
	\]
	积的切丛是两因子切丛拉回之和，因此 Whitney 公式给出
	\[
		w(\mathrm TX)=(1+a+a^2)(1+b+b^2).
	\]
	逐次展开，得到
	\[
		w_1=a+b,\quad w_2=a^2+ab+b^2,\quad
		w_3=a^2b+ab^2,\quad w_4=a^2b^2.
	\]
	于是 \(w_1^4=a^4+b^4=0\)，\(w_1^2w_2=(a^2+b^2)(a^2+ab+b^2)=2a^2b^2=0\)，且 \(w_1w_3=2a^2b^2=0\)。另一方面，\(w_2^2=a^2b^2=w_4\)，二者求值都为 \(1\)。

	对于 \(Y\)，复切丛给出 \(w_1=w_3=0\)、\(w_2=\overline h\)、\(w_4=\overline h^{\,2}\)，且 \(\langle\overline h^{\,2},[Y]_2\rangle=1\)。所以全部答案为
	\[
	\begin{array}{c|ccccc}
		& w_1^4&w_1^2w_2&w_1w_3&w_2^2&w_4\\ \hline
		X&0&0&0&1&1\\
		Y&0&0&0&1&1
	\end{array}
	\]
	这些数不能区分二者，且都证明 \(X\)、\(Y\) 各自不是无定向边界。在 \(X\amalg Y\) 上五个数相加后全部为零，所以这些数不再给出边界障碍。\cref{g2:thm:characteristic-number-bordism}证明的是配边推出各数相等；要由此表反过来断言二者配边，还需要这个必要条件的充分性，或另行构造具体配边。
\end{solution}

\begin{exercise}\label{g2:ex:connected-sum-bordism}
	设 \(M,N\) 为连通闭定向光滑 \(n\) 维流形，\(n\ge2\)。构造 \(M\amalg N\) 到定向连通和 \(M\mathbin{\#}N\) 的配边，并推出每个 Pontryagin 数在定向连通和下可加。判断 \(\mathbb CP^2\mathbin{\#}\mathbb CP^2\) 和 \(\mathbb CP^2\mathbin{\#}(-\mathbb CP^2)\) 是否为定向边界，并分别给出理由。
\end{exercise}
\begin{solution}
	从 \((M\amalg N)\times[0,1]\) 出发，在两个出边界上各选一张嵌入的 \(n\) 维闭球，沿这两张球附着一只 \(n+1\) 维的一把手 \([-1,1]\times D^n\)：它的两个端盘分别接到两因子，按反转边界方向的方式粘合。原来的上端边界被去掉两张球的内部，再接上把手侧面 \([-1,1]\times S^{n-1}\)。这正是定向连通和的颈部，故新的上端为 \(M\mathbin{\#}N\)。在球边缘的领圈坐标中，接缝是二维象限乘 \(S^{n-1}\)；把象限的折角换成一条与两侧相切的光滑弧，就得到光滑边界，并在该小邻域以外保持原构造。所得紧定向流形 \(V\) 满足
	\[
		\partial V=(-M)\amalg(-N)\amalg(M\mathbin{\#}N).
	\]
	因此 \([M\mathbin{\#}N]=[M]+[N]\) 于 \(\Omega_n^{\mathrm{SO}}\)。若 \(n=4k\)，每个 Pontryagin 数是这个群上的同态，故在连通和下等于两因子数值之和。

	\(\mathbb CP^2\mathbin{\#}\mathbb CP^2\) 的 \(p_1\) 数为 \(3+3=6\)，所以不是定向边界。\(\mathbb CP^2\mathbin{\#}(-\mathbb CP^2)\) 则有配边类 \([\mathbb CP^2]+[-\mathbb CP^2]=0\)，因而确是定向边界。这里边界存在性来自上述配边与反向圆柱的粘合；其 Pontryagin 数 \(3-3=0\) 与此相符。
\end{solution}

\begin{exercise}\label{g2:ex:pt-change-target-frame}
	设 \(f:S^{n+k}\to S^k\) 光滑，在基点附近恒取基点，且 \(0\) 为正则值。给定 \(A\in\operatorname{GL}(k,\mathbb R)\)，将 \(v\mapsto Av\) 延拓为一点紧化的自映射 \(\widehat A:S^k\to S^k\)。求 \(\widehat A\circ f\) 在零处的法标架，并分别说明 \(\det A>0\)、\(\det A<0\) 对应的稳定配边类。
\end{exercise}
\begin{solution}
	零点集合仍为 \(M=f^{-1}(0)\)。在法方向上，链式法则给出 \(\mathrm d(\widehat A\circ f)=A\,\mathrm df\)，故新标架为
	\[
		b_{\widehat A\circ f}=b_fA^{-1}.
	\]
	当 \(\det A>0\) 时，\(A^{-1}\) 可在 \(\operatorname{GL}^+(k,\mathbb R)\) 中连接到恒等，因此新旧标架同伦。若 \(\det A<0\)，可将 \(A^{-1}\) 连到反射 \(R=\operatorname{diag}(-1,1,\ldots,1)\)，于是稳定切标架变为 \(\phi_{b_f}(v,Ra)\)，即在一个平凡方向上反射。按逆元约定，这给出原稳定带标架配边类的负元。\(\widehat A\) 在无穷远处只需连续；复合映射在基点邻域恒定，在零点邻域光滑，若需要全局光滑代表，可保持这些区域作相对逼近。
\end{solution}

\begin{exercise}\label{g2:ex:cp-line-normal-framing-obstruction}
	把 \(\mathbb CP^1\) 作为标准射影直线嵌入 \(\mathbb CP^2\)。求其复法线丛的第一陈类，判断它能否成为某张光滑映射 \(\mathbb CP^2\to S^2\) 的正则值逆像。进一步判断增加平凡实方向能否使这张实法丛平凡。
\end{exercise}
\begin{solution}
	记 \(h\in H^2(\mathbb CP^1;\mathbb Z)\) 为正生成元。切丛限制分解为 \(\mathrm T\mathbb CP^2|_{\mathbb CP^1}\cong\mathrm T\mathbb CP^1\oplus\nu\)。由式\eqref{g2:eq:complex-projective-tangent-classes}，两边的第一陈类满足
	\[
		3h=2h+c_1(\nu),\qquad c_1(\nu)=h.
	\]
	法丛为复线丛，其底层定向实平面丛的 Euler 类就是 \(h\)，故不是平凡丛。如果它是 \(f^{-1}(y)\) 且 \(y\) 为正则值，微分会给出 \(\nu\cong\mathbb CP^1\times\mathrm T_yS^2\)，与此矛盾。

	稳定问题须另用稳定类判断。第21章的复丛实化比较给出 \(w_2(\nu_{\mathbb R})=c_1(\nu)\bmod2\ne0\)，而增加平凡实丛保持 \(w_2\)。所以它即使稳定后也不能平凡。Euler 类非零本身不足以完成最后一步，例如 \(\mathrm TS^2\) 的 Euler 类非零，却有 \(\mathrm TS^2\oplus\varepsilon^1\cong\varepsilon^3\)。
\end{solution}

\begin{exercise}\label{g2:ex:compact-level-framed-boundary}
	设 \(f:\mathbb R^{n+1}\to\mathbb R\) 光滑，\(0\) 为正则值，且在某个球外 \(f>0\)。证明 \(M=f^{-1}(0)\)，配上法向量 \(\nabla f\)，代表零元 \(\Omega_n^{\mathrm{fr}}\)。说明紧性假设在此有什么作用。
\end{exercise}
\begin{solution}
	子水平集 \(W=\{f\le0\}\) 闭且有界，因而紧。正则值定理给出 \(\partial W=M\)，外法向可取 \(\nu=\nabla f\)。环境标准切标架限制到 \(W\)，诱导的边界标架为
	\[
		(v,t)\longmapsto v+(-1)^nt\nabla f.
	\]
	因此这一稳定切标架代表零。法标架 \(b=\nabla f\) 所给的稳定切标架为 \((v,t)\mapsto v+t\nabla f\)：\(n\) 偶时它与上述标架相同，\(n\) 奇时它是上述类的逆元，仍为零。空水平集也满足结论。

	球外的正性使所取填充 \(W\) 紧，从而成为定义中的配边。若只知道 \(M\) 紧，负子水平集仍可能无界，例如 \(f(x)=1-|x|^2\)；此时不能直接把 \(\{f\le0\}\) 称作紧配边，须另选填充并核对标架。
\end{solution}

\begin{exercise}\label{g2:ex:standard-sphere-surgery}
	设 \(0\le p\le n-2\)，把 \(S^n\) 写成磨圆后的 \(\partial(D^{p+1}\times D^{n-p})\)。沿其中的 \(S^p\times D^{n-p}\) 使用乘积标架作手术，证明结果微分同胚于 \(S^{p+1}\times S^{n-p-1}\)，并核对 Euler 示性数的改变量。
\end{exercise}
\begin{solution}
	乘积圆盘磨圆后是 \(D^{n+1}\)，故其边界为标准 \(S^n\)。取走第一个面 \(S^p\times\operatorname{int}D^{n-p}\)，余下的是第二个面 \(D^{p+1}\times S^{n-p-1}\)。手术再接上一个这样的面，粘合映射在边界两个因子上都是恒等，于是结果是
	\[
		(D^{p+1}\cup_{S^p}D^{p+1})\times S^{n-p-1}
		\cong S^{p+1}\times S^{n-p-1}.
	\]
	这里两个圆盘的双倍采用通常的光滑领圈，磨圆得到标准球面。Euler 示性数之差为
	\[
		(1+(-1)^{p+1})(1+(-1)^{n-p-1})-(1+(-1)^n)
		=(-1)^{n-p-1}-(-1)^p,
	\]
	与\cref{g2:prop:surgery-euler}一致。特别地，\(S^4\) 沿标准圆周的手术给出 \(S^2\times S^2\)，其 Euler 示性数由 \(2\) 增至 \(4\)。
\end{solution}

\begin{exercise}\label{g2:ex:kill-fundamental-group-by-surgery}
	证明每个连通闭定向光滑 \(n\) 维流形 \(M\)，\(n\ge4\)，都与某个单连通闭定向光滑 \(n\) 维流形有向配边。请说明有限性、圆周嵌入、法标架及所得基本群四个环节。
\end{exercise}
\begin{solution}
	紧流形具有有限 CW 同伦型，故 \(\pi_1(M)\) 有有限生成集 \(g_1,\ldots,g_r\)。将代表这些元素的环路作光滑逼近，再对有限个圆周的并使用一般位置。当 \(n\ge4\) 时，圆周自身或两圆周之间的双点预期维数为 \(2-n<0\)，而微分为零的情形也可由一般位置避开。因此可用互不相交的嵌入圆周代表这些自由同伦类；基点道路的改变只产生共轭，不影响它们生成的正规子群。

	定向 \(M\) 和各圆周后，法丛获得定向。每个定向秩 \(n-1\) 实向量丛在圆周上平凡：沿割开圆周后的区间选标架，两端的差为 \(\operatorname{GL}^+(n-1,\mathbb R)\) 中一个矩阵；这个群道路连通，沿区间用一条道路修正标架，即能使两端吻合。管状邻域定理遂给出互不相交的 \(S^1\times D^{n-1}\) 邻域。

	对这些圆周逐个作定向兼容的 \(1\) 维手术。由于 \(1\le n-3\)，\cref{g2:prop:surgery-fundamental-group}适用，最终基本群为
	\[
		\pi_1(M)/\langle\!\langle g_1,\ldots,g_r\rangle\!\rangle=\{1\}.
	\]
	每步保持连通，并由其定向手术迹提供有向配边。把有限个迹沿共同边界粘合，就得到从 \(M\) 到最终单连通流形的有向配边。
\end{solution}

\begin{exercise}\label{g2:ex:surgery-trace-not-h-cobordism}
	设 \(M^n\) 闭且连通，\(n\ge4\)，在 \(M\) 上作一次 \(1\) 维手术，迹记为 \(W\)。证明 \(H_2(W,M;\mathbb Z)\cong\mathbb Z\)，从而 \(W\) 不是 \(h\)-配边。若附着圆周在 \(M\) 中可缩，解释为什么两端基本群同构仍不能改变这个结论。
\end{exercise}
\begin{solution}
	由\cref{g2:lem:handle-core-retraction}，流形对 \((W,M)\) 相对同伦等价于 \((M\cup_\alpha D^2,M)\)。相对胞腔链复形仅在维数 \(2\) 有一个 \(\mathbb Z\)，其他维数为零，所以其微分全为零，得到
	\[
		H_j(W,M;\mathbb Z)=
		\begin{cases}\mathbb Z,&j=2,\\0,&j\ne2.\end{cases}
	\]
	如果 \(M\hookrightarrow W\) 是同伦等价，所有相对同调群都应为零，矛盾。因此这个迹不是 \(h\)-配边。

	附着圆周可缩时，它生成的正规子群平凡，\cref{g2:prop:surgery-fundamental-group}确实给出 \(\pi_1(M_\varphi)\cong\pi_1(M)\)。然而基本群只检测一维环路，而新附着的核心是二维胞腔，相对 \(H_2\) 仍然记录它。故两端基本群相同不能代替 \(h\)-配边所要求的包含映射同伦等价。
\end{solution}

\begin{exercise}\label{g2:ex:trivial-whitehead-group}
	利用整数矩阵的初等变换证明 \(\operatorname{Wh}(\{1\})=0\)，并据此从 \(s\)-配边定理推出单连通的 \(h\)-配边定理。
\end{exercise}
\begin{solution}
	此时群环是 \(\mathbb Z\)，Whitehead 群为 \(K_1(\mathbb Z)/\langle[-1]\rangle\)。任一可逆整数矩阵 \(A\) 的行列式为 \(\pm1\)。其第一列各整数的最大公因数为 \(1\)，否则该公因数也整除行列式。使用整数行加法、行交换和变号，Euclid 算法可把第一列变为 \((1,0,\ldots,0)^{\mathsf T}\)；随后列加法消去第一行的其余元素。对剩余方阵归纳，得到单位矩阵。

	行加法和列加法对应初等矩阵。行交换可由初等矩阵及一个对角 \(-1\) 实现，变号本身也是对角 \(-1\)。故 \(K_1(\mathbb Z)\) 中每个类都由 \([-1]\) 生成，再除掉它便得到零群。

	现在设 \((W;M_0,M_1)\) 是两端闭连通、\(\dim W\ge6\) 的光滑 \(h\)-配边，且 \(M_0\) 单连通。其扭子只能落在零群 \(\operatorname{Wh}(\{1\})\)，所以必为零。\Cref{g2:thm:s-cobordism}给出固定 \(M_0\) 的微分同胚 \(M_0\times[0,1]\cong W\)，即所需 \(h\)-配边结论。
\end{solution}

\begin{exercise}\label{g2:ex:atlas-change-not-exotic}
	在实线的底层拓扑空间上，分别取由全局坐标 \(x\mapsto x\) 和 \(x\mapsto x^3\) 生成的最大光滑图册 \(\mathcal A\)、\(\mathcal B\)。证明这两个最大图册不同，但两个光滑流形微分同胚。进一步判断底空间的恒等映射在两个方向上是否光滑。
\end{exercise}
\begin{solution}
	坐标变换 \(x\mapsto x^3\) 光滑，但其逆映射 \(y\mapsto y^{1/3}\) 在原点不可微。因此两张坐标不属于同一个光滑图册，所生成的最大图册不同。

	设 \(h\colon(\mathbb R,\mathcal B)\rightarrow(\mathbb R,\mathcal A)\) 为底空间映射 \(h(x)=x^3\)。源端使用坐标 \(\beta(x)=x^3\)，靶端使用标准坐标 \(\alpha(x)=x\)，则其坐标表达为 \(\alpha h\beta^{-1}(y)=y\)。逆映射在相反的坐标中也表达为恒等映射，故 \(h\) 是微分同胚。

	底空间的恒等映射 \(\operatorname{id}\colon(\mathbb R,\mathcal A)\rightarrow(\mathbb R,\mathcal B)\) 的坐标表达是 \(x\mapsto x^3\)，所以光滑；反方向的坐标表达是立方根，所以不光滑。这说明恒等映射失去微分同胚性质，并不能排除其他重新识别底空间的微分同胚。
\end{solution}

\begin{exercise}\label{g2:ex:homotopy-seven-sphere-classes}
	设 \(\Sigma\) 是与 \(S^7\) 同胚的闭光滑流形，并任选一个定向。证明它的全部正次数 Stiefel--Whitney 类、Pontryagin 类以及切丛 Euler 类均为零。说明为什么这些类不能区分 \(M_{2,-1}\) 与标准 \(S^7\)。
\end{exercise}
\begin{solution}
	同胚给出 \(H^q(\Sigma;\mathbb Z)=H^q(S^7;\mathbb Z)\) 及模二版本，因此 \(H^q(\Sigma;\mathbb Z/2)=0\) 对 \(1\le q\le6\) 成立。于是 \(w_q(\mathrm T\Sigma)=0\) 对这些 \(q\) 成立。\(q>7\) 时，Stiefel--Whitney 类由秩条件为零。

	剩下 \(w_7\)。由\cref{g2:cor:tangent-euler-characteristic}，
	\[
		\langle e_{\mathrm{Th}}(\mathrm T\Sigma),[\Sigma]\rangle
		=\chi(\Sigma)=1+(-1)^7=0.
	\]
	\(H^7(\Sigma;\mathbb Z)\cong\mathbb Z\)，且在定向基本类上的求值是同构，故 \(e_{\mathrm{Th}}(\mathrm T\Sigma)=0\)。再由\cref{g2:thm:top-class-euler}，\(w_7\) 是该类的模二约化，因而也为零。

	Pontryagin 类 \(p_i\) 位于 \(H^{4i}(\Sigma;\mathbb Z)\)。\(i=1\) 时目标群为零，\(i\ge2\) 时次数超过维数，仍为零。这些消失结论同时适用于标准 \(S^7\) 和\cref{g2:thm:milnor-exotic-seven-sphere}中的 \(M_{2,-1}\)，但后两者不微分同胚。因而这组示性类在此没有足够的信息完成光滑分类。
\end{solution}

\begin{exercise}\label{g2:ex:exotic-r-four-proper-coordinate}
	设 \(R\) 是异光滑四维欧氏空间。证明不存在固有的单射光滑浸入 \(F\colon R\rightarrow\mathbb R^4\)。这里固有指任意紧集的原像紧。
\end{exercise}
\begin{solution}
	若这样的 \(F\) 存在，源、靶维数相同，浸入条件说明每点的微分都是线性同构。逆函数定理给出局部光滑逆，所以 \(F(R)\) 是开集。

	固有性又保证 \(F(R)\) 是闭集。具体地，设 \(y_j=F(x_j)\) 收敛于 \(y\)。集合 \(C=\{y\}\cup\{y_j:j\ge1\}\) 紧，故 \(F^{-1}(C)\) 紧。取 \(x_j\) 的收敛子列，极限记作 \(x\)；连续性给出 \(F(x)=y\)，所以 \(y\in F(R)\)。欧氏空间可度量，逐列闭性即闭性。

	\(F(R)\) 非空且开闭，而 \(\mathbb R^4\) 连通，故 \(F\) 满射。结合单射性，它是双射；各点的局部光滑逆在交叠处由单射性一致，组成整体光滑逆。因此 \(F\) 是微分同胚，与 \(R\) 异光滑矛盾。
\end{solution}

\begin{exercise}\label{g2:ex:exotic-r-four-one-point}
	设 \(R\) 是异光滑四维欧氏空间，\(R^+=R\cup\{\infty\}\) 是其单点紧化。
	\begin{enumerate}
		\item 证明 \(R^+\) 在拓扑上同胚于 \(S^4\)。若 \(R\) 的光滑结构能延拓成 \(R^+\) 上的光滑流形结构，证明所得闭流形不与标准 \(S^4\) 微分同胚。
		\item 假设进一步存在紧集 \(K\subset R\)，没有任何光滑嵌入的 \(S^3\subset R\) 能把 \(K\) 与无穷远端分开。证明在这种情况下，\(R\) 的光滑结构根本不能延拓过 \(\infty\)。
	\end{enumerate}
\end{exercise}
\begin{solution}
	1. 取同胚 \(h\colon R\rightarrow\mathbb R^4\)。同胚及其逆都将紧集映为紧集，因而令 \(h(\infty)=\infty\)，就得到单点紧化之间的同胚 \(R^+\rightarrow(\mathbb R^4)^+\cong S^4\)。若某个延拓后的光滑流形与标准 \(S^4\) 微分同胚，限制这个微分同胚到 \(R=R^+\setminus\{\infty\}\)，得到 \(R\cong_{\mathrm{DIFF}}S^4\setminus\{p\}\)。绕球面的一个旋转把 \(p\) 移到北极，再作球极投影，就给出右端与标准 \(\mathbb R^4\) 的微分同胚，产生矛盾。

	2. 若有延拓，\(\infty\) 有光滑坐标邻域。由于 \(K\) 紧且不含 \(\infty\)，可在这个邻域内取一个以 \(\infty\) 为中心、与 \(K\) 不交的小闭坐标四维球 \(B\)。其边界是 \(R\) 内的光滑嵌入 \(S^3\)。\(K\) 位于 \(R^+\setminus B\)，而 \(\infty\) 位于 \(\operatorname{int}B\)；从 \(K\) 到无穷远的任何连续路径都要穿过 \(\partial B\)。这与假设的不可分隔性质矛盾。
\end{solution}
